
\subsubsection{2.1 Alignment Benchmark Design}

TruthfulQA \citep{Lin2022} measures whether language models produce factually correct answers to questions that humans frequently answer incorrectly, using multiple-choice accuracy (MC2) as the primary metric. BBQ \citep{Parrish2022} is a question-answering benchmark designed to detect social biases across nine demographic categories, measuring response accuracy in ambiguous versus unambiguous contexts. Both benchmarks represent alignment-relevant constructs that are conceptually distinct from general reasoning or knowledge recall.

MMLU \citep{Hendrycks2021} evaluates broad academic knowledge across 57 tasks and serves as a standard proxy for general LLM capability in leaderboard evaluations. It has become a de facto scale indicator in comparative evaluations of open-weight models.

The Llama 2 technical report \citep{Touvron2023} provides a prominent example of joint alignment benchmark reporting: TruthfulQA MC2, BBQ, and safety scores are presented for base and chat model variants, with within-model improvements attributed to RLHF fine-tuning. Our work extends this comparison to the cross-model population by examining correlation structure across N=296 diverse open-weight models.

\subsubsection{2.2 Benchmark Structure and Dimensionality}

BenchScope \citep{Sha2026} characterizes the effective dimensionality of benchmark suites, reporting approximately 1.7 latent axes for the Open LLM Leaderboard, implying that alignment benchmarks are not fully reducible to a single capability dimension. This is consistent with our observation that a positive residual partial correlation remains after MMLU control: if effective dimensionality exceeds 1, then some benchmarks measure constructs beyond scale.

PCA-based analyses of cross-model benchmark variance \citep{clawrxiv2603} find that TruthfulQA loads primarily on PC2 (23.4\% of variance), a component orthogonal to PC1 (general capability). This structural result is consistent with our partial correlation results: controlling for MMLU (a PC1 proxy) reduces the TruthfulQA $\times$ BBQ correlation substantially. Our method uses directed partial regression rather than PCA, enabling a hypothesis test specific to the MMLU-TruthfulQA-BBQ triplet.

\subsubsection{2.3 RLHF and Alignment Training Effects}

Reinforcement learning from human feedback (RLHF; Stiennon et al., 2020; Ouyang et al., 2022) is the primary post-training method used to improve model alignment with human preferences. InstructGPT \citep{Ouyang2022} demonstrates improvements on human preference evaluations following RLHF. Llama 2 \citep{Touvron2023} documents RLHF-driven improvements on TruthfulQA and BBQ within-model. Whether RLHF produces systematic cross-model co-movement in bias proxies is examined in our optional Tier 3 analysis; the null result there (sign test p = 0.686) provides no cross-model evidence for this mechanism under the ARC proxy, though the null result must be interpreted cautiously given the proxy limitation.

Prior runs of this research pipeline established a within-family RLHF improvement on TruthfulQA MC2 of +3.406 points (BCa CI [2.589, 4.212]) across 321 base/chat pairs, consistent with the established literature.

\subsubsection{2.4 Scale Confounding in LLM Evaluation}

Scale confounding in LLM evaluation arises because larger, more capable models tend to score higher on nearly all benchmarks simultaneously, inflating apparent pairwise correlations. This is analogous to classical confounding in causal inference, where a shared common cause (scale) induces correlation between two outcomes. Partial correlation as a method for controlling confounders is well-established in biostatistics and social science \citep{Liu2017}, but has not been systematically applied to alignment benchmark correlation analysis.

The pingouin library \citep{Vallat2018Pingouin} implements partial Spearman correlation via an inverse covariance approach, returning exact p-values. The Fisher z transformation \citep{Fisher1915Frequency} provides a test for differences between dependent correlations estimated on the same sample. The combined application --- partial Spearman plus Fisher z difference test --- as an alignment benchmark diagnostic is, to our knowledge, novel.

